\cleardoublepage

\section{相关工作}

围绕新视角合成与高保真场景重建，本节先简要回顾传统三维重建算法和同步定位与地图构建（Simultaneous localization and mapping，简称 SLAM）相关算法。
随后介绍神经渲染、辐射场、3DGS 及其面向广角相机的扩展，并从这些工作中归纳出本文所关注的鱼眼标定、广角渲染和 LiDAR 初始化问题。

\subsection{传统三维重建算法}

早期的多视角合成任务主要以在运动结构恢复的算法为基础 \cite{snavely2006photo}，提出了从多张图像中恢复相机位姿和稀疏三维结构，依次进行最简单的新视角的合成。再之后SfM的算法经过不断的发展，逐渐收敛到SIFT特征提取\cite{lowe2004sift}、特征匹配、几何验证、增量式光束平差法优化的管线 \cite{schonberger2016sfm},该管线具有很强的鲁棒性，对于大多数场景都能达到很高精度的相机位姿标定与高质量的稀疏场景重建。
在获得稀疏几何之后，多视图立体重建（Multi-View Stereo, MVS）进一步从多视角观测中恢复稠密表面，如PMVS\cite{PMVS}、CMP-MVS\cite{jancosek2011multi}。这类方法具有几何模型清晰、可解释性强、工程体系成熟、易于与相机标定结合等优点，至今仍然是许多重建系统的基础。
同时随着深度学习的快速发展，深度学习网络开始在稀疏重建管线被应用。如 SuperPoint \cite{detone2018superpoint} 通过学习式的方法进行关键点与描述子提取，在弱纹理和较大视角变化场景中展现出更强的鲁棒性；
LightGlue \cite{lindenberger2023lightglue} 则通过深度学习网络进行稀疏局部特征的匹配，在保持较高匹配精度的同时显著提升了稀疏特征匹配效率。学习式特征与匹配器的引入提升了 SfM 的匹配数量与质量，进一步提高了 SfM 的鲁棒性与精度。
更进一步的方法如LoFTR\cite{sun2021loftr}则抛弃了特征检测的步骤，通过从粗到细的匹配流程解决了传统基于关键点的方法在弱纹理和重复纹理的区域的困难。而Detector-Free SfM\cite{he2024detector}则更进一步，先进行粗粒度的重建后，再不断迭代精修，从而在弱纹理区域得到更高的重建质量。

但是，SfM 本质上仍属于稀疏重建路线，即使学习方法提高了稀疏重建的相机精度和点云密度，其输出仍然只是稀疏点云，难以直接支撑高质量的新视角合成。即使通过 MVS 算法得到稠密点云，点云表示本身的离散性仍会在渲染时引入明显伪影。
因此，传统几何重建为本文提供了相机标定和几何初始化的基础，但其本身不能解决鱼眼图像下的高保真实时渲染问题；本文需要在可靠标定和几何先验之上进一步构建适合 3DGS 训练的图像与点云表示。

\subsection{神经渲染与辐射场算法}

深度学习网络在早期就被用于弱表面的重建\cite{jancosek2011multi}；卷积神经网络（CNN）也被用于基于图像的渲染(image-based rendering, 即IBR)的混合权重的预测 \cite{hedman2018deep}。这些方法虽然使用了神经网络，但是仍然是基于传统的多视角显式几何；而在最终渲染中使用CNN的方法会出现闪烁伪影。

基于体积表示的三维重建最早在 Soft3D\cite{penner2017soft} 被提出，其将深度图转化为基于体素的体积表示，通过投票进行多视图深度融合。随后的基于深度学习的体积步进算法也被提出 \cite{henzler2019escaping}, 使用一个连续可微的密度场进行几何表示。
再然后，NeRF 使用多层感知机将场景建模为连续的隐式辐射场\cite{mildenhall2021nerf},再通过体积步进的方法进行采样渲染。这样的方法隐式的对图片中的光照、几何、纹理进行建模，在图像保真度方面取得了显著进展。

然而，原始的NeRF使用了过大的多层感知机，严重影响了渲染和训练的效率。为解决这个问题，Instant-NGP 使用哈希网格进行神经编码 \cite{muller2022instant}，通过插值建模连续的神经场；Plenoxels 利用显式稀疏体素进行神经编码，\cite{fridovich2022plenoxels}，并通过球谐函数进行光场拟合。
这些工作的表示天生在空间上具有连续性，相当于将空间连续的先验加入了神经网络，降低了所需的神经网络规模，提高了优化和渲染的效率。
这些工作虽然都取得了显著的成果，但是隐式的场景建模和体积步进式的渲染仍然不可避免的需要在渲染时进行逐像素的多次采样，大量的无物体空间的采样带来了计算上的浪费。同时，图片质量也被有限的采样次数给制约，最后只能通过采样次数的调整进行质量和速度的权衡。
为解决这个问题，Kerbl等人提出了三维高斯泼溅（3D Gaussian Splatting, 3DGS）\cite{kerbl2023}，通过显式的各向异性的高斯粒子对场景进行建模，在兼顾渲染质量的同时达到了较高的渲染效率。
总体而言，神经辐射场路线证明了连续场景表示在新视角合成中的有效性，但其逐射线采样和体渲染开销仍不利于本文关注的快速室内重建管线。因此，本文选择以显式、可实时渲染的 3DGS 作为基础表示，并重点处理其在鱼眼相机和 LiDAR 初始化条件下的适配问题。

\subsection{3DGS及其扩展}

\subsubsection{3DGS 的基础方法}

3DGS是基于点的场景表示在新视角合成方向上的重要突破。与 NeRF 通过隐式神经辐射场表示场景不同，3DGS 直接使用一组可学习的三维高斯粒子作为场景表示。
每个高斯通常由中心位置、旋转、尺度、透明度以及颜色属性共同确定，其中颜色通常用球谐系数来建模视角相关的纹理，从而能够对高光和反射进行表达。每个三维高斯粒子几何上可看作定义在三维空间上的实值函数，写为
\begin{equation}
    \label{equ:3dgs-density}
    \rho : \mathbb{R}^{3} \rightarrow \mathbb{R},
    \qquad
    \rho(\mathbf{x})=\exp\left(-\frac{1}{2}(\mathbf{x}-\mu)^{T}\Sigma^{-1}(\mathbf{x}-\mu)\right).
\end{equation}
其中，\(\mathbf{x}\in\mathbb{R}^{3}\) 表示三维空间中的任意一点，\(\mu\in\mathbb{R}^{3}\) 表示该高斯的中心位置，\(\Sigma\in\mathbb{R}^{3\times 3}\) 表示控制高斯形状、朝向和空间扩张范围的协方差矩阵。直观上看，\(\mu\) 决定了高斯位于空间中的哪个位置，而 \(\Sigma\) 则决定了高斯沿不同方向的拉伸程度以及主轴朝向。

为了保证协方差矩阵始终半正定并便于优化，3DGS 通常并不直接优化 \(\Sigma\)，而是用旋转和尺度对其进行参数化，即
\begin{equation}
    \label{equ:3dgs-covariance}
    \Sigma = R S S^{T} R^{T}.
\end{equation}
其中，\(R\) 表示高斯主轴方向对应的旋转矩阵，\(S=\mathrm{diag}(s_{x},s_{y},s_{z})\) 表示沿三个主轴方向的尺度矩阵。由于 \(S S^{T}=\mathrm{diag}(s_{x}^{2},s_{y}^{2},s_{z}^{2})\)，因此高斯三个主方向由旋转矩阵 \(R\) 决定，而高斯沿各主方向的大小则分别由 \(S\) 决定。这样的几何表达使高斯既可以表示近似球形的小粒子，也可以表示细长、扁平等具有明确方向性的局部结构，具有高度的灵活性。

3DGS 的核心在于其可微渲染机制，使得每个高斯的参数都能得以有效训练。
对于每个三维高斯，作者先将三维高斯投影为二维椭圆高斯，从而可以有效计算其对二维像素平面的贡献。为实现这一过程，作者参考 \cite{zwicker2002ewa} 的思路，先通过局部线性近似将高斯的三维协方差矩阵投影为二维协方差矩阵：
\begin{equation}
    \label{equ:3dgs-projection}
    \Sigma' = J W \Sigma W^{T} J^{T}
\end{equation}
其中，\(W \in \mathbb{R}^{3 \times 3}\) 表示从世界坐标系到相机坐标系的线性变换矩阵，通常对应外参中的旋转部分；\(J \in \mathbb{R}^{2 \times 3}\) 表示在高斯中心附近对投影函数求得的雅可比矩阵，用于将非线性的投影过程局部近似为仿射变换。此时，第 \(i\) 个高斯粒子在像平面上的二维分布可写为

\begin{equation}
    \label{equ:3dgs-screen-density}
    \rho_i(\mathbf{v})=\exp\left(-\frac{1}{2}(\mathbf{v}-\mathbf{v}_{\mu_i})^{T}{\Sigma_i'}^{-1}(\mathbf{v}-\mathbf{v}_{\mu_i})\right).
\end{equation}

其中，\(\mathbf{v}\in \mathbb{R}^{2}\) 表示像平面上的一点，\(\mathbf{v}_{\mu_i}\in \mathbb{R}^{2}\) 表示第 \(i\) 个高斯中心在图像平面上的投影，\(\Sigma_i'\) 则表示该高斯投影到二维后的协方差矩阵。

最后，由于场景是多个高斯的混合，3DGS 采用基于粒子的体渲染方式，并通过球谐函数建模视角相关的颜色。假设当前要渲染的相机射线为 \(\mathbf{r}\)，其方向为 \(\mathbf{d}\in\mathbb{R}^{3}\)。渲染时，先将高斯按照其中心深度进行排序，再按如下方式进行近似的数值积分：
\begin{equation}
    \label{equ:3dgs-render}
    c(\mathbf{r}, \mathbf{d}) = \sum_{i=1}^{N} SH_i(\mathbf{d}) \, \alpha_i \prod_{j=1}^{i-1} (1-\alpha_j)
\end{equation}

其中，\(N\) 表示对当前射线产生贡献的高斯数量，\(SH_i(\mathbf{d})\) 表示第 \(i\) 个高斯在视线方向 \(\mathbf{d}\) 上由球谐函数给出的颜色值，\(\alpha_i\in\mathbb{R}\) 表示该高斯在当前像素处参与混合的有效透明度，其计算公式为
\begin{equation}
    \label{equ:3dgs-alpha}
    \alpha_i = \tau_i \rho_i(\mathbf{v}).
\end{equation}
其中，\(\tau_i\) 表示第 \(i\) 个高斯自身的透明度参数，\(\rho_i(\mathbf{v})\) 表示该高斯在像平面位置 \(\mathbf{v}\in\mathbb{R}^{2}\) 处的二维高斯贡献值，两者相乘后得到该高斯对当前像素的实际贡献强度。

在初始化阶段，3DGS 依赖 SfM 输出的稀疏点云进行初始化，然后通过各视角重建图像与真实图像的 \(l_1\) 损失和图像结构相似度损失（SSIM）进行监督优化。
同时 3DGS 引入了自适应的增密与剪枝机制：对于投影梯度较大、误差较高的区域，通过克隆或分裂高斯提升局部表示能力；对于透明度极低、尺度不合理的高斯，则进行裁剪和删除，从而能够自动的根据场景复杂度进行高斯数量的调整。
总体而言，3DGS 是以显式高斯表示为基础，通过可微光栅化渲染和误差驱动的梯度下降进行优化，并通过自适应的控制机制自动进行高斯点的增加和删除。这使它能够同时兼顾新视角合成质量、训练速度和渲染速度。

\subsubsection{3DGS 的质量与效率优化}

围绕 3DGS 的后续研究主要聚焦于几类问题：
其一是大场景下的渲染代价不稳定，需要引入多级细节（Level of Detail, LoD）机制；
其二是原始 3DGS 在几何一致性、抗锯齿和漂浮伪影方面仍存在缺陷，需要进一步提升视觉与几何质量；
其三是训练速度仍受优化与光栅化开销制约，需要更高效的训练算法和渲染管线。因此，当前 3DGS 的优化研究已经从单纯提高图像指标，扩展为对可扩展性、质量和效率的系统改进。

在大场景建模方面，V3DG\cite{yang2025v3dg}、Octree-GS \cite{ren2024octree} 和 CLoD-GS \cite{cheng2025clodgs} 等工作尝试通过显式层级结构或连续细节层次机制，缓解原始 3DGS 在远景视角下仍需渲染大量高斯的问题。这类方法的共同目标，是让场景表示能够随着观察距离变化而调整参与渲染的高斯规模，从而在保留近景细节的同时降低远景计算负担，使 3DGS 逐渐具备面向大场景和多尺度浏览的能力。
在质量提升方面，Mip-Splatting \cite{yu2024mipsplatting}通过在三维空间和图像空间引入更合理的滤波策略，缓解了原始 3DGS 在不同观察尺度下的混叠问题 ；2DGS 将三维体高斯改为二维高斯面片，并在优化中引入表面一致性约束，有效提升了几何精度 \cite{huang2024twodgs}；StableGS 则聚焦于漂浮伪影问题，通过跨视角深度一致性约束和更稳定的透明度建模减少非真实表面上的高斯 \cite{wang2025stablegs}；Geometry-Grounded Gaussian Splatting 则更进一步，试图从表示层面为高斯建立更明确的几何解释，以支持更稳定的深度恢复与表面提取 \cite{zhang2026geometrygroundedgs}。
在训练效率方面，一类工作着眼于优化算法本身。例如，3DGS$^2$ \cite{lan20253dgs2}尝试用二阶收敛的算法，从数学层面加速高斯的训练 。另一类工作如 Speedy-Splat\cite{hanson2025speedy}则更偏向系统层面的管线加速，通过计算更精确的高斯像素覆盖来降低训练和渲染开销。
这些 3DGS 扩展工作主要改进表示质量、训练效率或多尺度渲染能力，但大多默认输入来自透视相机和 SfM 稀疏点云。对于本文的鱼眼移动采集场景，仍需要同时处理强畸变投影、长序列标定和 LiDAR 点云初始化这三个前置问题。

\subsubsection{鱼眼与广角相机适配}

尽管 3DGS 在透视相机条件下表现优异，但其核心渲染流程存在对投影变换的局部线性近似，而这种近似源于 EWA splatting \cite{zwicker2002ewa} 中对高斯投影到2D的过程进行了投影函数为线性函数的假设。在普通视场角条件下，这一近似的误差通常可以接受；但对于鱼眼、超广角乃至超过 \(180^\circ\) 视场角的成像，非线性畸变会显著改变三维高斯在像平面上的投影形状，此时如果仍然沿用原始管线进行渲染，往往会在边缘区域产生不可忽略的误差，需要设计额外机制处理该部分误差。
3DGUT 是这一方向中较具代表性的工作 \cite{wu20253dgut}。该方法的目标是在保留 3DGS 原有高效 splatting 框架的同时，使其能够支持大畸变相机。其关键思想是利用无迹变换，在三维高斯上采样若干关键点，再将这些点在任意给定的相机内参模型中进行非线性投影，并根据投影后的点集去估计二维椭圆。这样做的优点在于将相机模型与固定透视投影公式解耦，让三维高斯的投影不再依赖于特定的投影公式；但其本质上仍然在尝试将三维高斯投影为二维椭圆，当视场角很大、边缘畸变很强时，三维高斯投影结果无法使用二维椭圆进行拟合，就会带来较大的投影误差。
相比之下，3DGEER 则放弃了对二维椭圆的依赖，而是直接在射线空间中对三维高斯进行体积分渲染 \cite{huang20253dgeer}。这种做法抛弃了直接将高斯投影为特定二维形状的做法，直接计算每个高斯沿任意相机射线的贡献，在原理上避免了由投影近似带来的误差。3DGEER 由于逐射线计算高斯贡献，对极端广角和鱼眼成像更为精确，但代价是较高的计算开销。
因此，现有鱼眼 3DGS 方法要么仍依赖二维椭圆近似，要么需要放弃高效光栅化管线。本文采用虚拟透视图拆分的思路，将强畸变鱼眼观测转换为若干小视场角透视观测，从而在避免直接鱼眼投影近似的同时复用标准 3DGS 渲染器。

\subsection{同步定位与地图构建算法}

与离线三维重建相比，同步定位与地图构建（SLAM）更强调在运动过程中实时估计传感器位姿并构建环境地图。
对本文的室内移动采集任务而言，SLAM 的核心价值在于为后续重建提供连续、稳定的位姿先验，降低后续重建过程中全局优化的难度。对于相机数量较多、轨迹较长的场景，若完全依赖图像间匹配和增量式注册进行 SfM，会面临多次全局光束平差优化计算量过大、相机参数容易过拟合的问题。
通过 SLAM 算法引入较精确的相机外参先验，一方面能够让后续场景优化收敛到更准确的解；另一方面，相对准确的先验内外参使本文能够直接进行全局光束平差优化，减少逐步注册和多次光束平差带来的时间开销。

基于激光雷达的3D SLAM算法以 LOAM\cite{zhang2014loam} 所提出的算法结构为起点，主要包含特征提取，里程计，建图三个组件。特征提取一般是提取激光雷达结果中类似于边缘或平面部分的点作为特征，里程计则是对相邻的两帧之间进行之间的特征进行匹配，并依次进行相邻帧之间的特征匹配。最后，多个时间扫描的结果被建图组件注册融合为一个完整的地图。
后续的基于激光雷达的工作仍然保持和LOAM相似的结构，比如 Lego-LOAM\cite{shan2018lego}使用点云分割算法滤除噪声，从而优化特征提取部分；并且引入了回环优化单元减少长程扫描的累积误差。 LOAM-Livox则将LOAM应用于新兴的固态激光雷达。为了解决小视场角和非重复扫描的问题，LOAM-Livox的里程计直接通过将新扫描配准全局地图来获得。

后续的工作则在激光雷达的基础上引入了惯性测量单元（IMU），提高了复杂环境下的轨迹恢复的精度与鲁棒性。FAST-LIO\cite{xu2021fast}引入了形式化的反向传播算法，精确的考虑了扫描中每个点的采样时间，并通过由IMU测量数据驱动的严格运动学模型来补偿运动畸变。同时使用创新性的卡尔曼增益公式，显著降低了计算量。 
FAST-LIO2 \cite{xu2022fastlio2}则基于FAST-LIO，尝试抛弃特征提取模块，直接将原始激光雷达点配准到地图中，进一步提升了精度与鲁棒性。同时提出了增量k-d树，能够实现真正实时大尺度建图。

除轨迹本身外，多传感器系统还面临跨模态外参标定问题。对于鱼眼相机与激光雷达的联合重建而言，两者之间存在一个刚性变化，需要合理的对齐算法能够将 LiDAR 所提供的几何信息对齐到图像观测中，从而能够根据odometry的外参得到相机较精确的预标定外参。
像素级外参自标定方法通过边缘或结构对齐来估计 LiDAR 与相机之间的坐标关系\cite{yuan2021pixel}，为跨传感器几何约束提供了可行方案。SLAM 及相关标定方法虽然无法直接得到三维高斯训练所需要的高精度标定结果，但它们可以为基于图像的SfM管线提供可靠的位姿初值和几何先验，从而有效提高SfM的效率和精度。
因此，SLAM 和跨传感器标定在本文中并不是最终重建结果，而是用于降低鱼眼 SfM 优化难度的几何先验。本文进一步将该先验与学习式特征匹配、三角化和光束平差结合，以获得可用于 3DGS 训练的相机参数。

\cleardoublepage
\section{主要方法}
本文提出了一套基于鱼眼相机和激光雷达的三维高斯场景重建管线，能够以双鱼眼相机采集得到的视频和激光雷达点云作为输入数据，实现室内三维高斯场景重建。
该方法的整体流程如图~\ref{fig:method-overview} 所示，主要解决快速运动结构恢复、大畸变鱼眼相机 3DGS 重建和激光雷达点云高效初始化三个问题，分别对应以下三个组件：
\begin{enumerate}
    \item 基于先验位姿的快速 SfM 管线（见第~\ref{sec:method-fast-sfm}节）：根据激光雷达、IMU 数据和预标定的相机到雷达外参得到鱼眼外参先验，用于提升鱼眼相机数据内外参优化与场景注册的效率和质量。
    \item 基于虚拟透视图的鱼眼图像三维高斯重建模块（见第~\ref{sec:method-perspective-sampling}节）：通过自动估计鱼眼有效水平、竖直视场角和虚拟透视图分辨率，将鱼眼图像拆分为八张共享光心的小视场角虚拟透视图，在保留大视场角信息的同时兼容传统三维高斯模型训练管线。
    \item 基于激光雷达的自适应初始化模块（见第~\ref{sec:method-lidar-mesh}节）：提出一种基于 K-d tree 的点云下采样与尺度感知初始化方法，使扫描更充分的区域在下采样后倾向于保留更多初始化点，并降低直接使用完整雷达点云带来的初始化规模。
\end{enumerate}

\begin{figure}[htbp]
    \centering
    \includegraphics[width=\textwidth]{paper/method-overview.pdf}
    \caption{本文方法整体流程图}
    \label{fig:method-overview}
\end{figure}

\subsection{基于几何先验的快速SfM管线}
\label{sec:method-fast-sfm}

鱼眼相机具有更大的视场角，但同时也带来了更强的畸变和更复杂的内外参耦合关系。若直接沿用传统增量式 SfM 管线，
一方面需要在长轨迹视频上反复执行注册与全局光束平差，同时相机外参和畸变参数的共同优化会降低光束平差的收敛稳定性，并且鱼眼相机的大视场角会带来更多三角化点，导致全局光束平差的单次迭代效率较低，两者结合导致传统 SfM 管线在鱼眼相机上的运行速度较慢；
而且鱼眼内参、畸变参数与位姿参数之间的耦合也给优化带来挑战，容易使内外参的优化陷入局部最优。为此，本文将 LiDAR-IMU 里程计提供的连续位姿作为相机外参先验，并结合迭代式 SfM 管线获得稳定的优化结果。

首先，本文使用 ChArUco 标定板预先估计鱼眼相机的焦距、主点和畸变参数，并将这些内参作为后续优化的初始化，从而减少大视场角成像下内参与位姿的相互过拟合；同时为相机到 LiDAR 的预标定过程提供相机内参。
然后，通过 FAST-LIO2 \cite{xu2022fastlio2}，本文利用原始激光雷达点云和 IMU 数据估计传感器在世界坐标系下的连续轨迹以及融合后的稀疏点云地图。
随后，结合基于边缘或结构对齐的 LiDAR-相机外参标定结果 \cite{yuan2021pixel}，即可通过里程计预测的外参得到相机坐标系的外参。记第 \(t\) 帧 LiDAR 到世界坐标系的刚体变换为 \(\mathbf{T}_{l_t}^{w}\)，相机到 LiDAR 的预标定外参为 \(\mathbf{T}_{c}^{l}\)，则相机初始位姿可写为
\begin{equation}
    \label{equ:method-camera-prior}
    \mathbf{T}_{c_t}^{w,0} = \mathbf{T}_{l_t}^{w}\mathbf{T}_{c}^{l}.
\end{equation}
式中 \(\mathbf{T}_{c_t}^{w,0}\) 为第 \(t\) 帧鱼眼图像在世界坐标系下的初始位姿。

在得到初始内外参后，本文不再完全依赖 COLMAP 的逐帧增量注册，而是利用该初始化的内外参数按照如下管线进行相机内外参优化。
首先，本文使用 SuperPoint\cite{detone2018superpoint} 进行鱼眼图像上的特征提取。
随后，本文使用 LightGlue\cite{lindenberger2023lightglue} 进行图像间特征匹配。考虑到完整场景扫描通常包含大量待标定图像，若使用全量两两匹配策略，匹配图像对数会随图像数量快速增长，带来较高计算开销。
为解决这个问题，本文设计了混合匹配对生成策略。匹配对由四部分组成：
一是滑动窗口匹配，即令每张图与其之后的 \(w\) 张图进行匹配，因为时序相近的图像通常也具有空间邻近性，\(w\) 常取 4；
二是根据空间信息生成匹配对，即对每张图，利用其先验外参选取 \(k\) 个邻近图像进行匹配，以保证当数据采集路径存在回环时，能够在回环区域进行一致性优化，\(k\) 一般取 20；
三是根据语义信息生成匹配对，即对每张图，使用 NetVLAD\cite{arandjelovic2016netvlad} 提取全图特征，然后选取 \(m\) 个语义邻近图像生成匹配对，以覆盖空间不相近但语义相似的图像，\(m\) 一般取 40；
四是长程时序锚点匹配，即每 \(g\) 帧取一个锚点，并对锚点两两匹配以增强全局标定一致性，其中 \(g\) 一般取 20，最后生成匹配对数约为 \(P_{anchor}(n) = \frac{n^2}{2 g^2}\)。
虽然锚点匹配对数量仍然包含二次项，但二次项具有较小常数系数，因此图像数量增长时匹配对数量不会过快增加，同时也可以通过参数调整控制匹配规模。

完成上述特征提取与匹配后，本文进行三次迭代以优化内外参。每次迭代使用上一轮相机参数重新三角化，再执行光束平差。
第一次迭代主要目的是对外参进行微调。由于基于 IMU 和激光雷达的外参估计会随距离增加产生累积误差，该轮使用较为宽松的重投影误差，并在光束平差优化时固定内参，以降低内外参共同优化进入局部极小值的风险。
第二次迭代主要目的是用更精确的三角化点进行微调，由于上一轮迭代时三角化使用的是激光雷达估计的质量较低的外参，在几何验证时可能会错误筛去很多匹配点，故这次迭代是为了得到更多的三角化点再次进行优化以提高标定质量，同时这次会将相机的内外参进行共同的优化。
第三次迭代主要目的是筛去错误匹配并进一步微调。由于此时已经得到置信度较高的内外参，该轮在三角化时使用更严格的重投影误差阈值，以筛去精度较低的匹配点，然后再用剩余匹配点对内外参进行共同优化，从而得到质量更高的 SfM 结果。

相对于渐进式的增量 SfM，本文通过激光雷达和 IMU 数据计算先验位姿，降低逐图像注册带来的累积误差，使所有图像在三角化后可以直接进入全局光束平差优化。该设计减少了逐步注册和多轮全局优化的计算开销，并为后续三维高斯训练提供相机参数。

\subsection{基于虚拟透视图的鱼眼图像三维高斯重建模块}
\label{sec:method-perspective-sampling}

传统 3DGS 主要面向针孔透视相机设计。对于普通相机而言，三维点首先由相机外参变换到相机坐标系，再通过针孔模型投影到归一化像平面。设相机坐标系下的点为 \(\mathbf{x}_c=(x,y,z)^{T}\)，则其归一化坐标为
\begin{equation}
    \label{equ:method-pinhole-normalized}
    \mathbf{m}=(a,b)^{T}=\left(\frac{x}{z},\frac{y}{z}\right)^{T}.
\end{equation}
真实镜头并不严格满足理想针孔模型，因此常在归一化像平面上继续引入畸变模型。以常见径向畸变为例，其核心思想是根据半径 \(r=\sqrt{a^2+b^2}\) 对归一化坐标做非线性扭曲，再通过内参矩阵映射到像素平面。也就是说，畸变参数并不改变相机具有单一光心的透视成像本质，而是在射线已经落到归一化平面之后，对二维坐标施加一个与半径相关的像平面上的变形。

鱼眼相机可以看作这一思想在大视场角条件下的扩展。常用的 KB-8 鱼眼模型\cite{kannala2006generic} 先将三维点变换到相机坐标系：
\begin{equation}
    \mathbf{x}_c=
    \left[
    \begin{array}{c}
        x \\
        y \\
        z
    \end{array}
    \right]
    =
    R\mathbf{X}+T.
\end{equation}
随后进行针孔模型下的归一化，并用归一化平面半径计算入射光线与光轴之间的夹角：
\begin{equation}
    a=\frac{x}{z},
    \qquad
    b=\frac{y}{z},
    \qquad
    r=\sqrt{a^2+b^2},
    \qquad
    \theta=\arctan r.
\end{equation}
然后，KB-8 模型用角度多项式将 \(\theta\) 映射为畸变后的角度半径：
\begin{equation}
    r_\theta
    =
    \theta
    \left(1+k_1\theta^2+k_2\theta^4+k_3\theta^6+k_4\theta^8\right)
    =
    \theta+k_1\theta^3+k_2\theta^5+k_3\theta^7+k_4\theta^9.
\end{equation}
该模型保持归一化平面上的径向方向不变，只改变半径大小：
\begin{equation}
    x'=\frac{r_\theta}{r}a,
    \qquad
    y'=\frac{r_\theta}{r}b.
\end{equation}
最后通过相机内参映射到像素坐标：
\begin{equation}
    \label{equ:method-kb8-projection}
    \pi_{\mathrm{KB8}}(\mathbf{x}_c)
    =
    \left[
    \begin{array}{c}
        f_xx'+c_x \\
        f_yy'+c_y
    \end{array}
    \right].
\end{equation}
OpenCV 的 fisheye 模型以及大部分被使用的畸变模型都采用这样的建模方式，注意到其重要性质为每个像素仍然对应一条从相机光心出发的三维射线，畸变模型或者说整个内参模型只是改变了这条射线所应该映射到的像素坐标。

基于这一观察，一个直接的想法是沿用传统 3DGS 对畸变图像的处理方式：根据标定得到的畸变参数，遍历目标无畸变图像的每个像素，通过双线性插值在原始畸变图像进行采样得到其对应的颜色，再使用标准透视相机 3DGS 进行训练。
然而对于本文使用的大视场角的鱼眼相机，若尝试得到与鱼眼的视场角相同的大视场角透视图，会导致在图像边缘部分的近处物体出现严重的伪影。因为在鱼眼边缘部分，单位距离所包含的像素数量较低，而在大视角的透视效应下，近处物体的单位距离则会对应更多的像素，
从而使得边缘区域的少量鱼眼像素会被拉伸到大面积透视平面上。如图~\ref{fig:method-high-fov-artifact} 所示，这会导致角分辨率不均、插值模糊和有效信息丢失，从而使得训练时 3DGS 的质量下降。

\begin{figure}[htbp]
    \centering
    \includegraphics[width=0.82\textwidth]{paper/high_fov_image.png}
    \caption{单张大视场角透视化重采样产生的伪影示例。图像边缘区域会被强烈拉伸，出现明显的模糊和结构缺失。}
    \label{fig:method-high-fov-artifact}
\end{figure}

因此，本文不再尝试用一张透视图承载完整鱼眼信息，而是提出利用多张共享同一光心的小视场角虚拟透视图来承载鱼眼采集到的信息。
在拆分时，本文允许虚拟透视图具有不同的水平和垂直视场角，并通过自动搜索确定竖直视场角、水平视场角、虚拟视图朝向和输出图像分辨率。不失一般性，本文假设鱼眼图像的有效水平视场角不小于竖直视场角；若某一相机不满足该条件，可以先转置图像并交换水平、竖直方向后再执行同样的搜索流程。

首先，本文根据鱼眼图像的有效像素掩码估计可用竖直视场角。记鱼眼有效像素掩码为 \(M_f\)，当像素位于相机画幅内且不属于边缘黑边区域时 \(M_f(\mathbf{x})=1\)，否则为0。本文用 \(F_y\) 表示鱼眼有效成像区域的完整竖直视场角，因此其竖直边界到光轴的最大入射角为 \(F_y/2\)。对一个候选竖直视场角 \(F_y\)，其在三维单位球面上对应一条竖直方向的候选边界：
\begin{equation}
    \label{equ:method-vertical-fov-boundary}
    \mathcal{C}_y(F_y)
    =
    \left\{
    \left[
    \begin{array}{c}
        \sin (F_y/2) \cos \varphi \\
        \sin (F_y/2) \sin \varphi \\
        \cos (F_y/2)
    \end{array}
    \right]
    \middle| \varphi \in [0,2\pi)
    \right\}.
\end{equation}
若 \(\mathcal{C}_y(F_y)\) 上的采样点投影到鱼眼图像后均落在有效像素区域内，则说明当前候选视场角仍偏小，可以继续增大；否则说明候选视场角已经超过有效成像范围，需要减小。由于该可行性随 \(F_y\) 增大而单调变差，本文在候选范围内通过二分搜索得到竖直方向的完整有效视场角。

在得到竖直视场角后，本文进一步搜索水平方向的有效视场角。直接由图像长宽比推导水平视场角只能给出一个保守初值，因为真实鱼眼有效区域会受到黑边、镜头裁切和标定模型共同影响。为此，本文以当前竖直视场角为固定条件，对候选水平视场角 \(F_x\) 做二分搜索，并对每个候选值执行一次低分辨率拆图测试。具体而言，本文在搜索阶段构造四个对角方向的虚拟透视视图，其中对角视角使用 \(F_x/6\) 与 \(F_y/6\) 的相对旋转，并检查每个虚拟图像边界像素反投影到原始鱼眼图像后是否仍位于有效图像范围内。若边界采样落到图像外或无效区域，说明当前 \(F_x\) 过大；反之说明该候选水平视场角仍可增大。该四视图测试仅用于降低水平视场角搜索开销，最终训练仍使用八张虚拟透视图。

确定 \(F_x\) 与 \(F_y\) 后，本文将每张鱼眼图像拆分为八张虚拟透视图。记 \(R_{\mathrm{pitch}}(\theta)\) 和 \(R_{\mathrm{yaw}}(\theta)\) 分别表示绕虚拟相机水平轴和竖直轴旋转 \(\theta\)。八个虚拟相机由四个轴向视角和四个对角视角组成，其中轴向视角分别覆盖上、下、左、右方向，对角视角覆盖四个角落区域。本文采用如下旋转集合：
\begin{equation}
    \label{equ:method-eight-view-rotations}
    \mathcal{R}
    =
    \mathcal{R}_{\mathrm{axis}}
    \cup
    \mathcal{R}_{\mathrm{diag}},
\end{equation}
其中轴向视角集合为
\begin{equation}
    \label{equ:method-eight-view-axis}
    \begin{array}{rl}
    \mathcal{R}_{\mathrm{axis}}
    &=
    \left\{
    R_{\mathrm{pitch}}\!\left(a\frac{F_y}{4}\right)
    \mid a\in\{-1,1\}
    \right\}
    \\
    &\quad{}\cup
    \left\{
    R_{\mathrm{yaw}}\!\left(b\frac{F_x}{4}\right)
    \mid b\in\{-1,1\}
    \right\}.
    \end{array}
\end{equation}
对角视角集合为
\begin{equation}
    \label{equ:method-eight-view-diag}
    \mathcal{R}_{\mathrm{diag}}
    =
    \left\{
    R_{\mathrm{yaw}}\!\left(b\frac{F_x}{6}\right)
    R_{\mathrm{pitch}}\!\left(a\frac{F_y}{6}\right)
    \mid a,b\in\{-1,1\}
    \right\}.
\end{equation}
每个虚拟透视图使用相同的 PINHOLE 相机模型。由于 \(F_x,F_y\) 表示原始鱼眼图像的完整水平和竖直视场角，本文令单张虚拟透视图的水平和竖直视场角分别为 \(F_x/2\) 和 \(F_y/2\)。因此，在针孔相机焦距计算中使用的半视场角分别为 \(F_x/4\) 和 \(F_y/4\)。若输出图像大小为 \(W\times H\)，则虚拟相机内参为
\begin{equation}
    \label{equ:method-perspective-k-xy}
    K_p(W,H,F_x/2,F_y/2)
    =
    \left[
    \begin{array}{ccc}
        \frac{W}{2\tan(F_x/4)} & 0 & \frac{W-1}{2} \\
        0 & \frac{H}{2\tan(F_y/4)} & \frac{H-1}{2} \\
        0 & 0 & 1
    \end{array}
    \right].
\end{equation}
在实际实现中，本文固定输出图像高为待搜索变量 \(N\)，并按原始图像长宽比设置输出宽度 \(W_N=\mathrm{round}(N W_f/H_f)\)，其中 \(W_f,H_f\) 为原始鱼眼图像宽高。记第 \(k\) 个虚拟透视相机相对于鱼眼相机的旋转为 \(R_k\)，对于虚拟透视图中的像素 \(\mathbf{u}=(u,v)^{T}\)，先通过透视相机反投影得到单位射线方向：
\begin{equation}
    \label{equ:method-virtual-view-ray}
    \mathbf{d}_{p,k}(\mathbf{u})
    =
    R_k
    \frac{K_p^{-1}\tilde{\mathbf{u}}}
    {\left\|K_p^{-1}\tilde{\mathbf{u}}\right\|}.
\end{equation}
根据前面以KB-8为例的鱼眼相机模型的投影性质，鱼眼投影函数只由入射光线方向决定，而与该光线上具体三维点到相机光心的距离无关。也就是说，对于任意距离 \(\lambda>0\)，同一条射线上的点都满足
\begin{equation}
    \label{equ:method-fisheye-ray-invariance}
    \pi_f(\lambda\mathbf{d}_{p,k}(\mathbf{u}))
    =
    \pi_f(\mathbf{d}_{p,k}(\mathbf{u})).
\end{equation}
与此同时，该射线由虚拟透视图中像素 \(\mathbf{u}\) 反投影得到，因此这条射线上的所有点在第 \(k\) 张虚拟透视图中也会投影到同一个像素 \(\mathbf{u}\)。因此，对于虚拟透视图上的任意像素 \(\mathbf{u}\)，都可以先确定其对应的三维射线方向 \(\mathbf{d}_{p,k}(\mathbf{u})\)，再通过鱼眼投影函数得到原始鱼眼图像中的连续像素坐标
\begin{equation}
    \label{equ:method-fisheye-sample-coordinate}
    \mathbf{x}_f(\mathbf{u})
    =
    \pi_f\!\left(\mathbf{d}_{p,k}(\mathbf{u})\right).
\end{equation}
最后，在原始鱼眼图像上对该连续坐标进行双线性插值，即可得到虚拟透视图中对应像素的颜色：
\begin{equation}
    \label{equ:method-fisheye-resample}
    I_k(\mathbf{u})
    =
    \mathrm{Bilinear}\!\left(I_f,\mathbf{x}_f(\mathbf{u})\right),
\end{equation}
其中，\(\tilde{\mathbf{u}}\) 为齐次像素坐标，\(I_f\) 为原始鱼眼图像，\(I_k\) 为第 \(k\) 张虚拟透视图。

上述过程可以理解为从虚拟透视图到原始鱼眼图像的反向采样。先构造虚拟透视相机内参和一组相对旋转，再为每个虚拟像素计算其在鱼眼图像中的采样坐标。由于每个虚拟像素都会通过双线性插值向原始鱼眼图像中的邻近像素贡献权重，本文可以反过来统计原鱼眼图像中每个有效像素被虚拟透视图覆盖的贡献值。
记 \(B(\mathbf{q},\mathbf{x})\) 为连续坐标 \(\mathbf{x}\) 对鱼眼像素 \(\mathbf{q}\) 的双线性插值权重，则在输出高度为 \(N\) 时，鱼眼像素 \(\mathbf{q}\) 的贡献值为
\begin{equation}
    \label{equ:method-fisheye-contribution}
    C_N(\mathbf{q})
    =
    \sum_{R_k\in\mathcal{R}}
    \sum_{\mathbf{u}\in\Omega_N}
    B\!\left(\mathbf{q}, \pi_f(\mathbf{d}_{p,k}(\mathbf{u}))\right),
    \qquad \mathbf{q}\in\Omega_f.
\end{equation}
其中，\(\Omega_N\) 表示当前输出高度下的虚拟透视图像素域，\(\Omega_f\) 表示原始鱼眼图像中的有效像素域。本文进一步对虚拟透视图输出高度 \(N\) 进行二分搜索，并给定覆盖阈值 \(\eta_c\) 和低覆盖比例阈值 \(\eta_r\)，选择满足
\begin{equation}
    \label{equ:method-contribution-condition}
    \frac{
    \left|\left\{\mathbf{q}\in\Omega_f \middle| C_N(\mathbf{q})<\eta_c\right\}\right|
    }{
    |\Omega_f|
    }
    < \eta_r
\end{equation}
的最小分辨率。这样可以保证绝大多数鱼眼有效像素达到预设覆盖强度，同时避免因为少量边界像素或标定噪声导致输出分辨率被过度放大。
完整流程如算法~\ref{alg:method-auto-perspective-sampling} 所示。

\begin{algorithm}[htbp]
    \caption{自动确定鱼眼虚拟透视图拆分方案}
    \label{alg:method-auto-perspective-sampling}
    \begin{algorithmic}[1]
        \REQUIRE 鱼眼投影函数 \(\pi_f\)，鱼眼有效像素掩码 \(M_f\)，原图大小 \(W_f\times H_f\)，覆盖阈值 \(\eta_c,\eta_r\)，输出高度范围 \([N_{\min},N_{\max}]\)
        \ENSURE 完整竖直视场角 \(F_y\)，完整水平视场角 \(F_x\)，八视图旋转集合 \(\mathcal{R}\)，虚拟透视图大小 \(W_N\times N\)
        \STATE 若有效水平视场角小于竖直视场角，则转置图像并交换水平、竖直方向
        \STATE 在候选范围内二分搜索 \(F_y\)
        \WHILE{\(F_y\) 搜索区间宽度大于阈值}
            \STATE 在 \(\mathcal{C}_y(F_y)\) 上均匀采样三维方向并投影到鱼眼图像
            \IF{所有投影点都位于有效像素区域}
                \STATE 增大 \(F_y\) 的下界
            \ELSE
                \STATE 减小 \(F_y\) 的上界
            \ENDIF
        \ENDWHILE
        \STATE 以长宽比推导的水平视场角作为初始可行下界，二分搜索 \(F_x\)
        \WHILE{\(F_x\) 搜索区间宽度大于阈值}
            \STATE 用 \(F_x/6,F_y/6\) 的对角旋转和 \(F_x/2,F_y/2\) 的虚拟视场角构造低分辨率四个对角虚拟视图用于搜索
            \STATE 将每个虚拟视图边界像素反投影并采样到原始鱼眼图像
            \IF{所有边界采样点均位于有效图像范围}
                \STATE 增大 \(F_x\) 的下界
            \ELSE
                \STATE 减小 \(F_x\) 的上界
            \ENDIF
        \ENDWHILE
        \STATE 根据式~(\ref{equ:method-eight-view-rotations})、式~(\ref{equ:method-eight-view-axis}) 和式~(\ref{equ:method-eight-view-diag}) 构造八视图旋转集合 \(\mathcal{R}\)
        \STATE 根据 \(F_x/2,F_y/2\) 构造虚拟透视相机内参 \(K_p\)
        \STATE 在 \([N_{\min},N_{\max}]\) 中二分搜索分辨率 \(N\)
        \STATE 对每个候选 \(N\) 计算所有有效鱼眼像素的贡献值 \(C_N(\mathbf{q})\)
        \STATE 返回满足式~(\ref{equ:method-contribution-condition}) 的最小 \(N\)
    \end{algorithmic}
\end{algorithm}

在确定 \(\mathcal{R}\) 与 \(N\) 后，即可执行鱼眼图像到虚拟透视图的反向采样，流程如算法~\ref{alg:method-virtual-perspective-sampling} 所示。

\begin{algorithm}[htbp]
    \caption{鱼眼图像到虚拟透视图的反向采样}
    \label{alg:method-virtual-perspective-sampling}
    \begin{algorithmic}[1]
        \REQUIRE 原始鱼眼图像 \(I_f\)，鱼眼投影函数 \(\pi_f\)，虚拟透视内参 \(K_p\)，虚拟相机旋转集合 \(\mathcal{R}\)，输出像素域 \(\Omega_p\)
        \ENSURE 虚拟透视图集合 \(\{I_k\}_{k=1}^{K}\) 与有效掩码集合 \(\{M_k\}_{k=1}^{K}\)
        \FOR{每个 \(R_k\in\mathcal{R}\)}
            \FOR{每个像素 \(\mathbf{u}\in\Omega_p\)}
                \STATE \(\mathbf{d}\Leftarrow R_k K_p^{-1}\tilde{\mathbf{u}} / \left\|K_p^{-1}\tilde{\mathbf{u}}\right\|\)
                \STATE \(\mathbf{x}_f\Leftarrow \pi_f(\mathbf{d})\)
                \IF{\(\mathbf{x}_f\) 位于有效鱼眼像素区域}
                    \STATE \(I_k(\mathbf{u})\Leftarrow \mathrm{Bilinear}(I_f,\mathbf{x}_f)\)
                    \STATE \(M_k(\mathbf{u})\Leftarrow 1\)
                \ELSE
                    \STATE \(M_k(\mathbf{u})\Leftarrow 0\)
                \ENDIF
            \ENDFOR
        \ENDFOR
    \end{algorithmic}
\end{algorithm}

从 3DGS 的渲染近似看，虚拟透视图的作用并不是改变三维高斯本身，而是避免在强鱼眼畸变投影上直接使用 EWA Splatting \cite{zwicker2002ewa} 的局部线性近似。若直接将三维高斯投影到鱼眼图像，投影函数可写为
\begin{equation}
    \label{equ:method-direct-fisheye-projection}
    \mathbf{x}_f
    =
    \pi_f(\mathbf{x}).
\end{equation}
此时若直接调用原始 3DGS 管线，需要在高斯中心附近对 \(\pi_f(\mathbf{x})\) 进行一阶近似，并用其雅可比矩阵将三维协方差投影为二维协方差。然而鱼眼投影中包含强非线性的径向畸变项，特别是在图像边缘区域，\(\pi_f\) 的投影过程具有高度非线性的性质，此时投影函数的一阶近似容易带来较大的误差，会把真实的弯曲的投影结果近似为不准确的二维椭圆。

虚拟透视图算法可以看作将鱼眼投影拆解为一个标准透视投影和一个从无畸变图像到畸变图像的变形。记第 \(k\) 张虚拟透视图的投影函数为 \(\pi_{v,k}\)，从虚拟透视图像素到原始鱼眼图像像素的映射为 \(w_{f,k}\)，则对于被该虚拟透视图覆盖的视场区域，有
\begin{equation}
    \label{equ:method-fisheye-decomposition}
    \pi_f(\mathbf{x})
    =
    w_{f,k}\!\left(\pi_{v,k}(\mathbf{x})\right),
    \qquad
    w_{f,k}(\mathbf{u})
    =
    \pi_f\!\left(\mathbf{d}_{p,k}(\mathbf{u})\right).
\end{equation}
\(\mathbf{d}_{p,k}(\mathbf{u})\) 表示第 \(k\) 张虚拟透视图中像素 \(\mathbf{u}\) 所对应的三维单位射线方向。具体而言，先用虚拟透视相机内参 \(K_p\) 将像素 \(\mathbf{u}\) 反投影到虚拟相机坐标系中的射线，再通过相对旋转 \(R_k\) 将该方向变换到原始鱼眼相机坐标系下。由于虚拟透视图与鱼眼图像共享同一光心，\(\mathbf{d}_{p,k}(\mathbf{u})\) 同时也是原始鱼眼相机中与该虚拟像素对应的入射光线方向。
经过这一拆解后，3DGS 在训练和渲染时只需要对 \(\pi_{v,k}\) 进行局部一阶近似，而由于每张虚拟透视图只覆盖较小视场角，且\(\pi_{v,k}\) 是普通透视投影，其多项式展开后投影中的非线性项远小于带有强畸变完整鱼眼投影中的非线性项，此时渲染中的一阶近似误差处于可接受范围内。而强非线性的鱼眼畸变不再进入三维高斯到二维高斯的近似过程，而是作为二维图像重采样中的 \(w_{f,k}\) 被精确处理。

从另一个角度而言，虚拟透视图中的二维椭圆高斯可以理解为：三维高斯先通过稳定的透视投影 \(\pi_{v,k}\) 得到在透视图上的投影结果，再通过 \(w_{f,k}\) 与原始鱼眼图像中的观测颜色建立对应。这个过程没有改变相机光心和三维射线，也没有改变空间中高斯与场景几何的关系，只是将鱼眼图像上的观测重新排列到若干张小视场角透视图中。因此，在标定准确的理想条件下，使用虚拟透视图训练与直接使用鱼眼射线训练具有等价性；实际误差主要来自离散图像插值和标定误差。

由于所有虚拟透视图与原始鱼眼图共享同一光心，其外参可以由鱼眼相机位姿直接传播得到。本文在该部分统一采用 camera-to-world 位姿约定，即 \(\mathbf{T}_{c}^{w}\) 表示从相机坐标系到世界坐标系的刚体变换。记原始鱼眼相机位姿为 \(\mathbf{T}_{f_t}^{w}\)，第 \(k\) 个虚拟透视相机到原始鱼眼相机的相对旋转为 \(R_k\)，则第 \(t\) 帧对应的第 \(k\) 张虚拟透视图位姿可写为
\begin{equation}
    \label{equ:method-perspective-pose}
    \mathbf{T}_{p_{t,k}}^{w} = \mathbf{T}_{f_t}^{w}
    \left[
    \begin{array}{cc}
        R_k & \mathbf{0} \\
        \mathbf{0}^{T} & 1
    \end{array}
    \right].
\end{equation}
这一设计能够在不改变场景几何中心的前提下，把鱼眼观测转换为传统 3DGS 可直接处理的透视相机集合。

除此之外，在进行透视图拆分时，位于边缘的图像仍然会出现透视图上多个像素由原始图像的单个像素得到的情况。这意味着准确标定需要原始鱼眼图像达到亚像素级标定精度，并且鱼眼标定误差会在边缘虚拟视图中被放大。
而3DGS的重建对外参的误差极度敏感，在外参不准的情况下易出现重影。
因此，为了缓解标定误差带来的影响，本文在虚拟透视图集合上再次运行与鱼眼标定相同的管线，根据预提供的内外参进行特征匹配、三角化和光束平差，对虚拟相机位姿做二次精化，然后将精化后的虚拟透视图和相机参数送入标准 3DGS 训练管线。

与直接重写鱼眼 3DGS 光栅化器相比，基于虚拟透视图的方法主要具有两点优势。
第一，它避免了单张大视场角去畸变图带来的极端透视拉伸，用多张小视场角图像更均匀地覆盖鱼眼相机捕捉到的宽视场信息；
第二，它可以复用现有透视相机 3DGS 渲染与优化框架，从而避免逐射线渲染带来的计算负担，提升训练效率。

\subsection{基于激光雷达的自适应初始化模块}
\label{sec:method-lidar-mesh}

仅依赖 SfM 稀疏点云初始化的 3DGS 通过自适应的控制模块在许多场景下能够得到可用结果。
然而，在实际场景中仍然面临两方面的挑战，首先是对于弱纹理墙面，纹理的单一导致该部分的几何训练是一个高度欠定的问题，其多解性会导致大量的漂浮伪影的产生；
而对于高细节区域，进行长序列扫描时，受训练时间的限制，每张图被采样的频率较低，在有限高斯点数量训练的前提下，有限的高斯容量会错误的被低频部分大量占用，导致一些高频细节部分高斯没能正确被密集化。
激光雷达能够提供尺度准确的密集三维点云，本文利用雷达点为三维高斯提供更充分的几何初始化，从而改善优化初值。

但是，如果直接使用全部雷达点初始化，会导致场景中高斯点的数量过多，带来巨大的训练和渲染负担。
因此需要对点云进行下采样，而普通体素下采样会倾向于把不同区域的点云在空间上压成近似均匀分布，反而破坏数据采集时天然形成的密度差异。
实际采集过程中，操作者通常会对具有更多几何细节或关注度更高的区域进行更充分的观察，这些区域在融合后的雷达点云中往往也具有更高密度。
本文希望保留这一性质，使高斯初始化在高细节区域天然更密，从而让后续优化更容易在这些区域具有更强表达能力。

为此，本文提出基于 K-d tree 的点云下采样方法。给定经过离群点剔除后的雷达点云 \(P=\{\mathbf{p}_i\}_{i=1}^{N}\)，假设目标最大初始化高斯点数为 \(N_t\)，首先根据目标点数自动确定叶节点容量
\begin{equation}
    \label{equ:method-kdtree-nmin}
    n_{\mathrm{leaf}}
    =
    \max\left(1,\left\lceil \frac{N}{N_t}\right\rceil\right).
\end{equation}
随后参考 K-d tree 的建立过程，递归地沿点云包围盒最长轴进行空间划分。对于某个待分裂节点 \(S\)，若其点数不超过 \(n_{\mathrm{leaf}}\)，则将其作为叶节点；否则在包围盒最长轴的中心位置进行划分。若中心划分导致某一侧为空，则退化为沿该轴排序后的中位数划分。对于每个叶节点 \(L_j\)，本文从其中随机选择一个原始雷达点作为代表点保留，并使用固定随机种子保证结果可复现：
\begin{equation}
    \label{equ:method-kdtree-centroid}
    \mathbf{p}_j^\ast
    =
    \mathrm{Sample}(L_j).
\end{equation}
如果最终叶节点数超过 \(N_t\)，则再随机裁剪到目标点数以内。该过程并不保证任意局部区域的点数比例严格不变，也不把点云密度保持作为理论约束；它只利用 K-d tree 的自适应空间划分性质，使扫描更充分、原始点数更多的区域通常被划分为更多局部叶节点，从而在下采样后倾向于保留更多代表点。由于实验中 LiDAR 初始化对整体平均指标的提升并不明显，本文将其定位为一种几何先验和初始化规模控制策略，而不是保证重建质量提升的核心模块。

\begin{algorithm}[htbp]
    \caption{基于 K-d tree 叶节点代表点的雷达点云下采样}
    \label{alg:method-kdtree-downsample}
    \begin{algorithmic}[1]
        \REQUIRE 雷达点云 \(P=\{\mathbf{p}_i\}_{i=1}^{N}\)，目标保留点数 \(N_t\)
        \ENSURE 下采样点云 \(P'\)，局部协方差集合 \(\{\Sigma_j\}\)
        \STATE \(n_{\mathrm{leaf}}\Leftarrow \max(1,\lceil N/N_t\rceil)\)
        \STATE 初始化节点队列 \(\mathcal{Q}\Leftarrow\{(P,P)\}\)，叶节点及父节点集合 \(\mathcal{L}\Leftarrow\emptyset\)
        \WHILE{\(\mathcal{Q}\) 非空}
            \STATE 从 \(\mathcal{Q}\) 中取出一个节点及其父节点 \((S,A)\)
            \IF{\(|S|\leq n_{\mathrm{leaf}}\)}
                \STATE 将 \((S,A)\) 加入集合 \(\mathcal{L}\)
            \ELSE
                \STATE 选择 \(S\) 的包围盒中长度最大的坐标轴 \(a\)
                \STATE 按第 \(a\) 维包围盒中心将 \(S\) 划分为 \(S_l\) 和 \(S_r\)
                \IF{\(S_l\) 或 \(S_r\) 为空}
                    \STATE 沿第 \(a\) 维排序，并按中位数重新划分 \(S_l,S_r\)
                \ENDIF
                \STATE 将 \((S_l,S),(S_r,S)\) 加入 \(\mathcal{Q}\)
            \ENDIF
        \ENDWHILE
        \STATE \(P'\Leftarrow\emptyset,\ \{\Sigma_j\}\Leftarrow\emptyset\)
        \FOR{每个叶节点及其父节点 \((L_j,A_j)\)}
            \STATE 从 \(L_j\) 中随机选择代表点 \(\mathbf{p}_j^\ast\)
            \STATE 计算父节点 \(A_j\) 中点坐标的协方差矩阵 \(\Sigma_j\)
            \STATE 将 \(\mathbf{p}_j^\ast\) 加入 \(P'\)，将 \(\Sigma_j\) 加入协方差集合
        \ENDFOR
        \IF{\(|P'|>N_t\)}
            \STATE 随机裁剪 \(P'\) 和对应协方差至 \(N_t\) 个元素
        \ENDIF
        \STATE 返回 \(P'\) 和 \(\{\Sigma_j\}\)
    \end{algorithmic}
\end{algorithm}

完成下采样后，本文将保留的雷达点转换为 3DGS 可读取的高斯初始化文件。对于第 \(j\) 个叶节点，代表点 \(\mathbf{p}_j^\ast\) 直接作为高斯中心 \(\mathbf{\mu}_j\)。若原始点云包含颜色属性，则将 RGB 颜色归一化到 \([0,1]\)，并写入球谐函数的直流分量
\begin{equation}
    \label{equ:method-lidar-sh-dc}
    \mathbf{f}_{\mathrm{dc},j}
    =
    \frac{\mathbf{rgb}_j-0.5}{C_0},
    \quad
    C_0=0.2820947918,
\end{equation}
其余高阶球谐系数初始化为 0；若点云没有颜色属性，则使用白色作为初始颜色。透明度参数对所有高斯统一初始化为 1。

高斯的初始尺度和旋转由代表点所在叶节点的父节点协方差给出。相比直接使用叶节点内的点统计协方差，父节点包含更大的局部邻域，能够更稳定地估计该区域点云的表面分布。记第 \(j\) 个代表点所在叶节点的父节点为 \(A_j\)，其协方差矩阵为
\begin{equation}
    \label{equ:method-lidar-covariance}
    \Sigma_j
    =
    \frac{1}{|A_j|-1}
    \sum_{\mathbf{p}\in A_j}
    (\mathbf{p}-\bar{\mathbf{p}}_{A_j})(\mathbf{p}-\bar{\mathbf{p}}_{A_j})^{T},
\end{equation}
其中 \(\bar{\mathbf{p}}_{A_j}\) 为父节点内点坐标均值。当局部点云协方差退化时，本文使用最小尺度 \(s_{\min}\) 进行正则化。对 \(\Sigma_j\) 进行特征分解并按特征值从大到小排序：
\begin{equation}
    \label{equ:method-lidar-eigen}
    \Sigma_j = U_j \mathrm{diag}(\lambda_{j,1},\lambda_{j,2},\lambda_{j,3}) U_j^{T}.
\end{equation}
若 \(U_j\) 的行列式为负，则翻转最后一列以得到合法旋转矩阵。随后将 \(U_j\) 转换为四元数作为高斯旋转参数，并将尺度初始化为
\begin{equation}
    \label{equ:method-lidar-scale}
    \mathbf{s}_j
    =
    \log\left(
    \gamma
    \sqrt{
    \max(\lambda_j, s_{\min}^{2})
    }
    \right),
\end{equation}
其中 \(\gamma\) 为尺度放大系数，实验中取 3。局部点云的协方差矩阵能够反映该区域的表面性质：在平面或边缘附近，较大的特征值通常对应表面切向或结构延展方向，较小的特征值对应法向或厚度方向。因此，用父节点协方差初始化高斯的尺度和旋转，可以使初始各向异性高斯更贴合局部表面几何，而不是从各向同性点状高斯开始优化。

\cleardoublepage
\section{实验设计与结果分析}

本章围绕本文提出的鱼眼相机与激光雷达融合三维高斯重建管线进行实验验证。实验主要回答三个问题：
第一，基于虚拟透视图的鱼眼图像重建是否能够在保留大视场角信息的同时避免单张大视场角去畸变图带来的边缘拉伸伪影；
第二，几何先验辅助的 SfM 管线是否能够提升鱼眼图像内外参优化的效率与质量；
第三，基于激光雷达的自适应初始化与 K-d tree 下采样是否能够在控制初始化规模的同时提供可用的几何先验。

\subsection{评价指标}

为了验证新视角合成质量，本文使用该方向较为常用的指标 PSNR、SSIM\cite{wang2004image} 和 LPIPS\cite{zhang2018unreasonable} 来衡量渲染质量，从像素误差、结构相似度、感知相似度三个角度对新视角合成结果进行定量分析。

\subsection{数据集介绍}
本章通过多个数据集评估本文所提出的重建管线的有效性。
在具体实验中，本文首先使用 ScanNet++ 的验证场景进行主要实验。该数据集的 DSLR 目录中提供鱼眼原图、相机参数以及其扫描点云。
对于 Zip-NeRF，本文使用其公开的鱼眼图像与对应的相机参数。
训练/测试划分采用 \(llffhold=8\) 的间隔采样策略：按采集顺序每相邻 8 张鱼眼图像中抽取 1 张作为测试源图像，其余 7 张作为训练源图像。训练阶段只使用训练源图像生成的虚拟透视图；测试阶段不使用测试源图像参与优化，而是以测试鱼眼图像中心方向为光轴，按照固定 FoV 拆出一张小视场角透视图作为评价视图。所有方法在相同测试透视图上计算 PSNR、SSIM 和 LPIPS，以保证不同相机输入域和不同标定结果之间的对比一致。

本文第一部分量化实验在公开通用畸变相机数据上测试算法表现，测试场景来自 ScanNet++\cite{yeshwanth2023scannet++} 和 Zip-NeRF。
其中，ScanNet++ 提供室内场景的鱼眼图像、去畸变图像与扫描点云，可用于验证本文的鱼眼虚拟透视图构造和雷达点初始化模块；Zip-NeRF 提供鱼眼图像及其去畸变版本，可用于验证不同相机输入域下的新视角合成质量。
所有参与该部分量化实验的数据使用的相机 FoV 均为 \(180^\circ\)。

\subsection{实现细节}

为保证对比公平，所有方法使用相同的训练/测试划分、图像分辨率、训练迭代次数和高斯数量上限，各方法按照其公开实现中的默认设置进行数据训练。
对于本文方法，本文采用基于 Easyvolcap\cite{xu2023easyvolcap} 优化后的训练工具用于 3DGS 的训练，训练的每个迭代会抽取单张鱼眼训练图，用其对应的虚拟透视图进行梯度累计后再单次反向传播进行训练。
在训练方面，本文采用3DGS-MCMC的密集化方案，在 scannet++ 数据集上限定高斯数上限为 2M，Zip-NeRF数据集上限定高斯数上限为5M；不透明度正则化的权重设为0；尺度正则化的权重设为0.0001。
所有实验均使用单个 RTX 4090 D GPU 和 251GB RAM 作为默认配置，如有修改将在对应实验中特别说明。


\subsection{各数据集上的量化对比}

表~\ref{tab:exp-main-fisheye} 为各数据集上的主要量化结果。Scene 列同时标注场景名称及其数据集来源；Method、PSNR、SSIM 和 LPIPS 列在同一单元格中按相同顺序列出四种对比方法及对应结果，便于逐场景填写和比较。
在这一章节中，各方法均使用鱼眼相机图像作为训练输入，并使用由测试鱼眼图像中心方向去畸变得到的小视场角透视图作为测试图像。由于重建后的三维场景在实际应用中通常会以透视相机视角进行渲染，而不是重新渲染为鱼眼图像，因此该设置用于验证不同算法从鱼眼训练输入迁移到透视测试视角的新视角合成能力。

\begingroup
\scriptsize
\setlength{\tabcolsep}{4pt}
\begin{longtable}{@{}
    >{\raggedright\arraybackslash}p{0.18\textwidth}
    >{\raggedright\arraybackslash}p{0.30\textwidth}
    >{\centering\arraybackslash}p{0.12\textwidth}
    >{\centering\arraybackslash}p{0.12\textwidth}
    >{\centering\arraybackslash}p{0.12\textwidth}
    @{}}
    \caption{各测试场景上的新视角合成质量量化结果}
    \label{tab:exp-main-fisheye}\\
    \toprule
    Scene & Method & PSNR$\uparrow$ & SSIM$\uparrow$ & LPIPS$\downarrow$ \\
    \midrule
    \endfirsthead
    \caption[]{各测试场景上的新视角合成质量量化结果（续）}\\
    \toprule
    Scene & Method & PSNR$\uparrow$ & SSIM$\uparrow$ & LPIPS$\downarrow$ \\
    \midrule
    \endhead
    \bottomrule
    \endfoot
        \makecell[l]{apartment\\ScanNet++} &
        \makecell[l]{3DGS + 去畸变图像\\3DGUT\\3DGEER\\本文方法} &
        \makecell[c]{28.85\\24.40\\31.37\\32.66} &
        \makecell[c]{0.958\\0.898\\0.963\\0.974} &
        \makecell[c]{0.080\\0.274\\0.075\\0.066} \\
        \midrule
        \makecell[l]{bedroom\\ScanNet++} &
        \makecell[l]{3DGS + 去畸变图像\\3DGUT\\3DGEER\\本文方法} &
        \makecell[c]{30.55\\25.38\\32.354\\32.34} &
        \makecell[c]{0.971\\0.908\\0.971\\0.980} &
        \makecell[c]{0.067\\0.278\\0.066\\0.062} \\
        \midrule
        \makecell[l]{classroom\\ScanNet++} &
        \makecell[l]{3DGS + 去畸变图像\\3DGUT\\3DGEER\\本文方法} &
        \makecell[c]{30.016\\24.59\\33.41\\33.84} &
        \makecell[c]{0.954\\0.884\\0.961\\0.976} &
        \makecell[c]{0.098\\0.321\\0.094\\0.074} \\
        \midrule
        \makecell[l]{kitchen\\ScanNet++} &
        \makecell[l]{3DGS + 去畸变图像\\3DGUT\\3DGEER\\本文方法} &
        \makecell[c]{30.465\\22.96\\30.767\\32.77} &
        \makecell[c]{0.956\\0.848\\0.953\\0.968} &
        \makecell[c]{0.096\\0.317\\0.090\\0.085} \\
        \midrule
        \makecell[l]{office\\ScanNet++} &
        \makecell[l]{3DGS + 去畸变图像\\3DGUT\\3DGEER\\本文方法} &
        \makecell[c]{27.33\\24.37\\28.489\\30.17} &
        \makecell[c]{0.913\\0.873\\0.925\\0.937} &
        \makecell[c]{0.156\\0.309\\0.146\\0.129} \\
        \midrule
        \makecell[l]{alameda\\Zip-NeRF} &
        \makecell[l]{3DGS + 去畸变图像\\3DGUT\\3DGEER\\本文方法} &
        \makecell[c]{20.54\\22.108\\22.274\\22.41} &
        \makecell[c]{0.665\\0.673\\0.686\\0.703} &
        \makecell[c]{0.430\\0.523\\0.437\\0.400} \\
        \midrule
        \makecell[l]{berlin\\Zip-NeRF} &
        \makecell[l]{3DGS + 去畸变图像\\3DGUT\\3DGEER\\本文方法} &
        \makecell[c]{25.87\\26.45\\26.49\\27.56} &
        \makecell[c]{0.881\\0.877\\0.887\\0.904} &
        \makecell[c]{0.264\\0.384\\0.278\\0.238} \\
        \midrule
        \makecell[l]{london\\Zip-NeRF} &
        \makecell[l]{3DGS + 去畸变图像\\3DGUT\\3DGEER\\本文方法} &
        \makecell[c]{26.09\\25.878\\26.13\\27.04} &
        \makecell[c]{0.793\\0.792\\0.802\\0.815} &
        \makecell[c]{0.414\\0.499\\0.430\\0.395} \\
        \midrule
        \makecell[l]{nyc\\Zip-NeRF} &
        \makecell[l]{3DGS + 去畸变图像\\3DGUT\\3DGEER\\本文方法} &
        \makecell[c]{26.505\\26.573\\26.94\\27.61} &
        \makecell[c]{0.797\\0.791\\0.794\\0.815} &
        \makecell[c]{0.367\\0.460\\0.452\\0.350} \\
\end{longtable}
\endgroup

从表~\ref{tab:exp-main-fisheye} 可以看出，本文方法在 ScanNet++ 和 Zip-NeRF 两类公开鱼眼场景上均取得了稳定的量化优势。在 ScanNet++ 室内场景中，本文方法在 apartment、classroom、kitchen 和 office 等场景上取得最高的 PSNR，并在所有 ScanNet++ 场景上取得最高的 SSIM 和最低的 LPIPS；即使在 bedroom 场景中 PSNR 与最优方法接近，本文方法仍在结构相似度和感知误差上保持优势。这说明基于虚拟透视图的训练方式不仅能够缓解鱼眼边缘的强畸变问题，也能在室内遮挡和复杂几何结构下保持较好的整体重建质量。

在 Zip-NeRF 场景上，本文方法同样在 alameda、berlin、london 和 nyc 四个场景中取得最优或接近最优的结果，并在感知指标 LPIPS 上表现出更稳定的优势。与直接在去畸变图像上训练 3DGS 相比，本文方法避免了单张大视场角去畸变图带来的边缘拉伸；与 3DGUT 和 3DGEER 等直接处理畸变相机或射线空间的方法相比，本文方法通过多个小视场角虚拟透视图复用标准 3DGS 的透视投影假设，一方面使得训练的模型可以复用当前很多已有的3DGS的渲染器进行渲染， 另一方面消除了鱼眼的畸变对训练的影响而没有增加更多的近似，提高了训练结果的质量。
因此，主实验结果表明，本文方法的提升并非局限于某一特定数据集，而是能够在不同来源、不同场景布局的鱼眼数据上稳定成立。

\begin{figure}[htbp]
    \centering
    \fbox{
    \begin{minipage}[c][0.34\textheight][c]{0.94\textwidth}
        \centering
        \small 定性网格对比图占位：按场景展示 GT、3DGS + 去畸变图像、3DGUT、3DGEER 和本文方法的渲染结果，并在鱼眼边缘、弱纹理墙面和高频细节区域放大局部误差。
    \end{minipage}
    }
    \caption{各数据集上的新视角合成定性对比图}
    \label{fig:exp-qual-grid}
\end{figure}

图~\ref{fig:exp-qual-grid} 将用于补充表~\ref{tab:exp-main-fisheye} 的定量结论。定量指标能够反映整体重建误差，而通过上述真值图片和不同方法重建结果的渲染图对比，更加直接的展示了我们的方法对重建质量的提升，可以看到我们的方法可以更好的重建出场景中的细节纹理。



\subsection{SfM 管线实验}

为了验证几何先验辅助的 SfM 管线在真实采集长序列数据上的效率与精度，本文在真实采集的 company 和 jazz 两个场景上比较传统增量式 SfM 标定结果与本文快速 SfM 管线得到的相机参数。两组相机参数分别送入相同的虚拟透视图生成流程和相同的 3DGS 训练设置，最终使用测试视角上的 PSNR、SSIM 和 LPIPS 衡量由相机参数支撑的新视角合成质量。同时，本文记录两种 SfM 管线在相同硬件环境下完成标定所需的运行用时，用于比较增量式 SfM 与本文快速 SfM 的效率差异。

本文不将重投影误差作为主要评价指标。原因在于重投影误差主要反映已匹配特征点在当前优化模型下的一致性，容易受到特征分布、外点过滤和可见点数量影响；对于 3DGS 训练而言，更关键的是相机外参是否能在稠密图像监督下具有良好的几何一致性，从而能支撑稳定的新视角合成。因此，本实验采用渲染质量作为间接但更贴近下游任务的评价标准，并结合 SfM 标定用时共同评价两条管线的效果。

\begin{table}[htbp]
    \centering
    \small
    \setlength{\tabcolsep}{4pt}
    \caption{增量式 SfM 与本文快速 SfM 管线的标定用时和新视角合成质量对比}
    \label{tab:exp-sfm}
    \begin{tabularx}{\textwidth}{
        >{\raggedright\arraybackslash}p{0.18\textwidth}
        >{\raggedright\arraybackslash}p{0.24\textwidth}
        >{\centering\arraybackslash}X
        >{\centering\arraybackslash}X
        >{\centering\arraybackslash}X
        >{\centering\arraybackslash}X}
        \toprule
        Scene & 标定方法 & SfM 用时$\downarrow$ & PSNR$\uparrow$ & SSIM$\uparrow$ & LPIPS$\downarrow$ \\
        \midrule
        company & 增量式 SfM & 待填 & 待填 & 待填 & 待填 \\
        company & 本文快速 SfM & 待填 & 待填 & 待填 & 待填 \\
        \midrule
        jazz & 增量式 SfM & 待填 & 待填 & 待填 & 待填 \\
        jazz & 本文快速 SfM & 待填 & 待填 & 待填 & 待填 \\
        \bottomrule
    \end{tabularx}
\end{table}

表~\ref{tab:exp-sfm} 的结果用于验证本文快速 SfM 管线在真实采集数据上的实用价值。实验结果表明，本文快速 SfM 管线在显著缩短标定用时的同时，取得了不低于传统增量式 SfM 的新视角合成质量。该结果说明，LiDAR-IMU 提供的先验位姿能够减少长序列鱼眼图像中的逐帧注册开销，并且也不会因为累积误差导致后续的三角化于光束平差的优化进入局部极小值。

从下游重建角度看，快速 SfM 的意义不仅在于节省标定时间，也在于得到足够稳定的相机外参。 而 3DGS 对外参误差非常敏感，轻微的姿态偏差就可能在训练中表现为重影、浮点或局部纹理错位。因此，company 和 jazz 场景中的结果说明，先验位姿辅助的迭代式标定管线并没有以牺牲重建质量为代价换取速度，而是同时提高了真实采集长序列数据的处理效率和工程可用性。

\subsection{ScanNet++ 上的消融实验}

本文在 ScanNet++ 上通过消融实验分析各模块对最终重建结果的贡献。
实验从最基础的去畸变 3DGS 开始，依次加入虚拟透视图、虚拟透视图位姿二次精化和雷达点初始化对结果的影响，每个部分消融后的量化结果见表~\ref{tab:exp-ablation}。

\begin{table}[htbp]
    \centering
    \small
    \caption{本文各模块的消融实验}
    \label{tab:exp-ablation}
    \begin{tabularx}{\textwidth}{
        >{\raggedright\arraybackslash}p{0.46\textwidth}
        >{\centering\arraybackslash}X
        >{\centering\arraybackslash}X
        >{\centering\arraybackslash}X}
        \toprule
        方法设置 & PSNR$\uparrow$ & SSIM$\uparrow$ & LPIPS$\downarrow$\\
        \midrule
        去畸变图像训练 & 29.44 & 0.950 & 0.327 \\
        虚拟透视图 & 30.998 & 0.961 & 0.094 \\
        +位姿二次精化& 31.97 & 0.966 & 0.085 \\
        +K-d tree 下采样后三维高斯初始化 & 32.36 & 0.967 & 0.083 \\
        \bottomrule
    \end{tabularx}
\end{table}

消融结果首先表明，虚拟透视图是本文方法中贡献最大的模块。与直接使用去畸变图像训练相比，加入虚拟透视图后 PSNR 和 SSIM 明显提升，LPIPS 大幅下降。这一结果与方法设计一致：虚拟透视图将单张大视场角鱼眼图像拆分为多个小视场角透视图，从而避免在强非线性鱼眼投影上直接使用 3DGS 的局部线性投影近似，也减少了大视场角去畸变图在边缘区域产生的拉伸伪影。

位姿二次精化进一步提高了虚拟透视图之间的几何一致性。由于鱼眼边缘区域的像素在虚拟透视图重采样后会被放大，原始标定中的微小误差也会在透视图边缘表现为更明显的错位。消融实验中，加入位姿二次精化后各项指标继续改善，说明在虚拟透视图集合上再次进行特征匹配、三角化和光束平差，能够缓解鱼眼标定误差向下游 3DGS 训练传播的问题。

K-d tree 下采样后的 LiDAR 高斯初始化带来了较小但稳定的指标变化。该现象说明，在训练迭代充分且 3DGS 具备自适应密集化机制时，最终平均重建质量主要仍受相机标定精度、图像覆盖和优化过程影响，额外 LiDAR 点云先验并不必然带来显著的全局指标提升。因此，本文不将 LiDAR 初始化表述为保证质量提升的核心因素，而将其定位为一种尺度准确的几何先验和初始化规模控制策略：它可以为弱纹理区域提供更合理的初始几何分布，也能通过父节点协方差初始化得到更贴合局部表面结构的各向异性高斯。


\subsection{结果分析和局限性}

综合主实验、SfM 管线实验和消融实验可以看出，本文方法的主要收益来自标定质量提升和投影模型适配两个方面。
主量化实验表明，基于虚拟透视图的训练方式能够在 ScanNet++ 和 Zip-NeRF 上稳定优于去畸变训练、通用畸变相机 3DGS 方法和直接射线空间渲染方法；消融实验进一步说明，该优势主要来自虚拟透视图对鱼眼强畸变投影的分解，以及位姿二次精化对多虚拟视图几何一致性的修正。这些结果共同支撑了本文的核心观点：对于大视场角鱼眼场景，与其直接修改 3DGS 光栅化器以适配复杂畸变投影，不如将鱼眼观测转换为一组标定一致的小视场角透视观测后，直接复用成熟的透视相机 3DGS 训练框架进行训练。

快速 SfM 管线实验进一步验证了本文方法在真实采集流程中的工程价值。通过快速 SfM 我们能够让 SfM 的用时维持线性增长，而不是像传统增量式管线中，图像数量的增加会导致每次全局光束平差的耗时和全局光束平差的次数都大大上升，从而导致超线性的时间增长。
从实际工程的角度上考虑，标定时间的降低使得我们能够以更高的fps从视频中提取图像，从而降低两帧之间的视差，使得我们可以进一步提高标定中点的轨迹长度，从而带来更好的标定质量；同时，更密集的图像提取有利于对高光和反射区域的重建，因为密集的观测能够对这样的区域具有更加良好的约束，带来更好的重建效果。

然而，本文方法仍存在一定局限性。首先，虚拟透视图的质量依赖原始鱼眼内外参和虚拟视图位姿的准确性，当标定误差较大时，重采样过程可能在边缘区域放大像素错位，因此仍需要位姿二次精化来缓解这一问题。
其次，我们的快速 SfM 管线实际上对初始化具有较为严格的精度要求，我们当前初始化利用的是实时 SLAM 的结果，其本身就已经有较高的精度，每张图与精确的标定结果的误差大概在十个像素以内，在此基础上快速SfM管线才能得到比较好的效果。
然后，LiDAR 初始化在平均 PSNR、SSIM 和 LPIPS 上提升有限，说明在训练充分的设置下，3DGS 的自适应控制机制已经能够重新分配高斯点，雷达先验更多是让弱纹理区域有比较好的初始化，使得低纹理区域产生更少的浮点。
最后，本文实验主要覆盖室内和公开鱼眼场景，未来仍需要在更大尺度的场景中进一步验证该管线的鲁棒性。

\section{总结与展望}

本文围绕鱼眼相机与激光雷达融合条件下的三维高斯场景重建问题，提出了一套面向室内大视场角采集数据的完整重建管线。针对鱼眼图像中存在的强畸变、长序列标定开销大以及弱纹理区域几何初始化不足等问题，本文没有直接对 3DGS 光栅化器进行复杂改写，而是从相机标定、图像表示和几何初始化三个层面进行处理，使鱼眼采集数据能够更稳定地接入标准三维高斯训练框架。
首先，本文提出了基于几何先验的快速 SfM 管线。该管线利用 LiDAR-IMU 里程计和相机到雷达的预标定外参，为鱼眼图像提供较准确的相机位姿初值，并在此基础上进行图像匹配、三角化和光束平差优化。与传统增量式 SfM 相比，该策略减少了逐帧注册和反复全局优化带来的计算开销，使长序列鱼眼数据的标定过程更加高效。同时，真实采集场景上的实验表明，该管线能够在缩短标定用时的同时保持比较优秀的标定结果，为后续高质量三维高斯的重建提供较为良好的位姿。
同时，本文提出了基于虚拟透视图的鱼眼图像三维高斯重建方法。该方法通过估计鱼眼图像的有效视场角和输出分辨率，将单张鱼眼图像拆分为多张共享光心的小视场角虚拟透视图，并将每张虚拟透视图作为标准透视相机输入 3DGS 训练。这样既保留了鱼眼相机的大视场角观测信息，又避免了单张大视场角去畸变图像在边缘区域产生严重拉伸，同时减少了直接在鱼眼投影下使用局部线性高斯投影近似所带来的误差。实验结果表明，该模块有效提升了基于鱼眼相机重建的高斯的质量。
其次，本文引入了基于 LiDAR 点云的高斯初始化策略。针对完整雷达点云规模过大、直接初始化会增加训练负担的问题，本文采用 K-d tree 下采样控制初始高斯数量，并利用代表点父节点的局部协方差初始化高斯尺度和旋转，使初始化的高斯能够更好适配局部表面几何。消融实验表明，在训练充分的情况下，LiDAR 初始化也能起到一定的提升效果，更加重要的是，它为弱纹理区域和局部表面结构提供更合理的初始分布。
最后，本文在 ScanNet++、Zip-NeRF 以及真实采集场景上对所提出的方法进行了验证。各数据集上的量化对比表明，本文方法相较去畸变图像训练、通用畸变相机 3DGS 方法和直接射线空间渲染方法，在 PSNR、SSIM 和 LPIPS 等指标上取得了更稳定的表现；消融实验进一步验证了虚拟透视图、位姿二次精化和 LiDAR 初始化各模块的作用。总体而言，本文证明了在鱼眼和激光雷达联合采集条件下，通过几何先验辅助标定、虚拟透视图重采样和尺度感知初始化，可以构建一条兼顾标定效率、渲染质量和工程可用性的三维高斯重建管线。

本文方法对 LiDAR 融合得到的密集点云仍有进一步利用空间。实际实验中，基于融合后的雷达点云，经过下采样并结合 GP3 算法 \cite{marton2009fast} 已经可以得到较为完整的场景网格；然而在真实采集场景中，网格与相机位姿之间仍存在精确对齐的困难，使得本文尚未充分利用网格法线和深度信息对三维高斯的位置与协方差施加几何约束。若后续工作能够进一步解决网格、点云与相机观测之间的高精度配准问题，则有望将更强的表面几何先验引入三维高斯优化过程，从而提升模型的几何一致性和跨视角泛化能力。

此外，本文当前的标定管线虽然能够取得较好的结果，但其精度仍然依赖后续的虚拟透视图位姿二次优化。实验分析表明，限制标定精度的主要因素在于鱼眼图像上的特征提取和匹配质量。对于基于学习的 SuperPoint 特征，由于其主要在低畸变图像上训练，直接应用于鱼眼图像时特征点质量和匹配稳定性会下降，难以形成足够数量的有效匹配；对于传统 SIFT 特征，鱼眼图像中的强畸变会导致同一三维点在不同视角下的局部外观变化较大，同样不利于构建稳定的长匹配轨迹。
因此，未来工作可以进一步围绕鱼眼图像中的长轨迹匹配展开，探索利用 LiDAR 密集点云为跨视角图像匹配提供几何引导。通过将雷达几何先验与图像特征匹配过程结合，有望提升鱼眼相机标定的精度和鲁棒性，并进一步减少后续位姿二次精化对重建质量的影响。
